\runninghead{Where the Body Is}
\properlane{where-the-body-is}{Where the Body Is: The Scattered Gathered to the Crucified and Risen Christ}

Between the lightning of the coming and the darkening of the sun the Gospel
sets a proverb: \latin{Ubicúmque fúerit corpus, illic congregabúntur et
áquilæ}, ``Wheresoever the body shall be, there shall the eagles also be
gathered together'' (Mt 24:28). Three verses later the angels ``gather together
his elect from the four winds'' (v.~31). The proverb is about carrion birds and
a carcass, and the Fathers do not read it so. They agree that the body is
Christ's and the eagles are the saints; they differ on where the gathering
happens. This interpretation reads the end that the Gospel describes as an
assembly of the scattered to Christ crucified and risen, and finds that
assembly already begun in two of the Mass's texts. The Introit promises to
bring back the captivity ``from all places,'' and the Epistle thanks the
Father who has ``translated us into the kingdom of the Son of his love'':
each moves people from where they were to where God brings them. For the
Introit the interpretation follows Jerome, who reads the promise, beyond its
first sense, as the return of penitent sinners to the Church.

\subsection{The body that falls in death}
\label{sec:body-jerome}

St Jerome begins from the natural fact behind the proverb. Eagles, he says,
sense a carcass even across the sea and gather to it; how much more, he
argues, must we \latin{festinare ad eum, cujus fulgur exit ab Oriente},
``hasten to him whose lightning goes out from the east'' (\work{In
Matthaeum}~IV, on 24:28). He then reads the word itself. The Greek of the
Gospel has \latin{ptoma}, a fallen thing: \latin{Possumus autem corpus, id
est, ptoma, quod significantius Latine dicitur cadaver, ab eo quod per mortem
cadat, passionem Christi intelligere, ad quam provocamur; ut ubicumque in
Scripturis legitur, congregemur, et per illam venire possimus ad Verbum Dei},
``We can understand the body, that is, \latin{ptoma}, which in Latin is more
expressively called \latin{cadaver} because it falls by death, as the Passion
of Christ, to which we are summoned, so that wherever it is read in the
Scriptures we may gather, and through it may come to the Word of God.'' He
cites the psalm in which the sufferer says ``They have dug my hands
and feet'' and the prophecy of the one led ``as a sheep to the slaughter'' (Ps
21:17; Is 53:7).

The eagles in Jerome's account are \latin{sancti, quibus innovata est juventus
ut aquilae}, ``the saints, whose youth is renewed like the eagle's'' (Ps
102:5), and who, in Isaias's words, ``take wings as eagles'' (Is 40:31)
\latin{ut ad Christi veniant passionem}, ``so as to come to the Passion of
Christ.'' For Jerome the gathering happens wherever the Passion is read. The
word he chose, \latin{congregemur}, ``that we may gather,'' is the Gospel's
own verb for the eagles.

\subsection{The place of the Passion}
\label{sec:body-hilary}

St Hilary of Poitiers places the gathering in the place where Christ
suffered: \latin{Sanctos de volatu spiritalis corporis aquilas nuncupavit:
quorum, congregantibus angelis, conventum futurum in loco passionis ostendit.
Et digne illic claritatis adventus exspectabitur, ubi nobis gloriam
aeternitatis passione corporeae humilitatis operatus est}, ``He called the
saints eagles from the flight of the spiritual body, and showed that their
assembly, when the angels gather them, will be in the place of the Passion.
And it is fitting that the coming of his brightness be awaited there, where he
worked for us the glory of eternity by the passion of his bodily lowliness''
(\work{In Matthaeum}~25.8). Hilary's reason is the fitness of the place.

He reads the signs in the sky the same way. They show \latin{gloriam adventus
sui et claritatis reditum}, ``the glory of his coming and the return of his
brightness,'' and the angels sent with the trumpet to gather the saints. He
lists likenesses of lowly things made great, the great tree from the mustard
seed and the light set on the wood among them, and ends \latin{sic ex
humilitate mortis Dei gloria}, ``so, from the lowliness of death, the glory of
God.'' By all these signs, he says, the Lord willed us to know
\latin{tempus nostrae redemptionis, quo ex corruptione corporum in honorem
spiritalis substantiae transferemur}, ``the time of our redemption, in which we
shall be carried over from the corruption of bodies into the honour of a
spiritual substance'' (26.1). Hilary's verb for that last passage,
\latin{transferemur}, is the verb the Epistle uses in the past tense when it
thanks the Father who \latin{tránstulit}, ``hath translated'' us into the
kingdom of his Son (Col 1:13).

\subsection{The body taken into heaven}
\label{sec:body-augustine}

St Augustine reads the body as the Lord's risen body and the place as heaven:
\latin{Ubi fuerit corpus, illo congregabuntur aquilae; id est in coelum, quo
hinc secum levavit corpus in homine suscepto: quod etiam cadaver ideo
appellatum est, quia moriturus haec loquebatur}, ``Where the body is, thither
the eagles will be gathered: that is, into heaven, whither he raised with him
from here the body, in the manhood he had taken up; and it was called a
carcass because he was about to die when he said these things.'' The eagles
are \latin{spirituales, qui ejus passionem humilitatemque imitando, tanquam de
ejus corpore saturantur}, ``the spiritual, who, by imitating his Passion and
lowliness, are fed as it were from his body'' (\work{Quaestiones
Evangeliorum}~I.42). Augustine's \latin{tanquam}, ``as it were,'' keeps the
feeding a likeness: the spiritual are nourished by imitating the Passion.

St John Chrysostom adds the angels. The coming ``needs not a herald'' but shows
itself everywhere at once, and the eagles are ``the multitude of the angels, of
the martyrs, of all the saints'' (\work{Hom.\ in Mt.}~76.3). For him the place
is wherever Christ appears. He reads the sign of the Son of man as ``the cross
being brighter than the sun,'' and observes that the Lord comes ``no longer on
the cross, but `in the clouds of Heaven'\,'' (76.3). St Thomas Aquinas, explaining the
word \latin{cadaver}, says that it signifies the Passion \latin{quia tunc
veniet Christus ostendens signa passionis}, ``because Christ will then come
showing the signs of his Passion,'' and he reports Jerome in his own words:
\latin{ubicumque sit memoria Christi passionis, viri sancti debent congregari},
``wherever the memory of Christ's Passion is, holy men ought to gather''
(\work{Super Matthaeum}, cap.~24, citing Heb 10:32).

The four answers do not agree on the place: the Passion as read and
remembered, the place of the Passion, heaven, wherever Christ appears. They
agree that the gathering is to Christ, and Jerome, Hilary and Aquinas name the
Passion as what draws it.

\subsection{Brought back from all places}
\label{sec:body-introit}

The Introit's antiphon ends \latin{et redúcam captivitátem vestram de cunctis
locis}, ``and I will bring back your captivity from all places.'' Between
those two phrases Jeremias has \latin{et congregabo vos de universis gentibus},
``and I will gather you out of all nations'' (Jer 29:14), and the antiphon
does not sing it. The verb it passes over is the one the Gospel uses for the
eagles, \latin{congregabúntur}, and for the angels who \latin{congregábunt
eléctos eius}, ``shall gather together his elect,'' ``from the four winds,
from the farthest parts of the heavens to the utmost bounds of them'' (Mt
24:28, 31). Since God's promise to the exiles and the Lord's description of
the end use one word for the gathering of God's people, the return from all
places may be taken as a first instance of the gathering from the four
winds.

St Jerome, commenting on the letter itself, gives it a sense beyond the return
from Babylon. \latin{Secundum anagogen, tamdiu in saeculi hujus confusione
versamur, quamdiu requiem septenarii numeri mereamur accipere, et accepta
poenitentia, Deus impleat quod promisit, et reducamur ad locum nostrum
Ecclesiam}, ``According to the anagogy, we are tossed in the confusion of this
world until we are found worthy to receive the rest of the number seven, and,
our penance accepted, God fulfils what he promised, and we are brought back to
our own place, the Church'' (\work{In Hieremiam}~V, on 29:10--14, PL~24
col.~893). On the next verses he writes, \latin{et invocabimus eum, et
revertemur in Ecclesiam, et orabimus, et exaudiemur}, ``and we shall call on
him, and return into the Church, and pray, and be heard'' (col.~894). The
exiles of this sense are sinners, and the place to which they are brought back
after penance is the Church.

St Augustine asks which captivity the verse of the Introit means, since
\latin{puto enim neminem nostrum modo esse apud Barbaros}, ``I think that none
of us is now among the barbarians.'' He answers with Paul's ``law of sin''
that is ``captivating me'' (Rom 7:23), and says the psalm is spoken to Christ,
who turned away that captivity \latin{non liberando nos a Barbaris \dots\ sed
liberando nos ab operibus malis, a peccatis nostris, per quae nobis satanas
dominabatur}, ``not by freeing us from barbarians, but by freeing us from evil
works, from our sins, through which Satan ruled over us'' (\work{Enarrationes
in Psalmos}~84.4). The release is forgiveness: \latin{iniquitas tenebat
captivum, remissa iniquitate liberaris}, ``iniquity held you captive; when
iniquity is forgiven you are set free.'' He draws a duty from it at once:
\latin{Confitere ergo te esse in captivitate, ut dignus sis liberari: nam qui
hostem suum non intellexit, quomodo invocat liberatorem?}, ``Confess, then,
that you are in captivity, that you may be worthy to be set free: for he who
has not understood his enemy, how shall he call upon the deliverer?'' (84.5).

\subsection{Delivered and translated}
\label{sec:body-epistle}

The Epistle speaks of deliverance in the past tense and says where it
leads. The Father \latin{erípuit nos de potestáte tenebrárum, et tránstulit in
regnum Fílii dilectiónis suæ}, ``hath delivered us from the power of darkness
and hath translated us into the kingdom of the Son of his love'' (Col 1:13).
Chrysostom notices both verbs. Paul said ``delivered'' and not ``sent forth,''
``showing our great misery''; he said ``translated,'' as one leads a soldier
over from one position to another, ``so that it is both of us and of Him.''
Those ``that were enemies, those that were in darkness'' were brought ``to
where the Son is, to the same honor with Him'' (\work{Hom.\ in Col.}~2). St
Thomas Aquinas glosses the power of darkness as \latin{sive Daemonum, sive
peccatorum}, ``whether of demons or of sins,'' and quotes Isaias: ``even the
captivity shall be taken away from the strong'' (Is 49:25; \work{Super
Colossenses}, cap.~1, lect.~3). Paul's ``lot of the saints in light'' (Col
1:12) is the company into which the delivered are brought, and his ``In whom
we have redemption through his blood'' (v.~14) names the price.

The Gradual gives thanks for deliverance: \latin{Liberásti nos, Dómine, ex
affligéntibus nos}, ``Thou hast delivered us, O Lord, from them that afflict
us.'' Rupert of Deutz, a twelfth-century abbot who wrote on the liturgy of
the Church's year, reads these chants as the story of
Israel's remnant at the end.
In the Introit God promises a banquet of grace to his brethren according to
the flesh, as Joseph, ruling Egypt, received the brothers who had sold him and
feasted with them (Gen 43:16, 34). The remnant will come back to the one
\latin{quem negaverunt et occiderunt}, ``whom they denied and killed,'' and
\latin{ab eo recipientur in gratiam}, ``will be received by him into grace.''
\latin{Soluti \dots\ captivitate animarum}, ``freed from the captivity of
souls,'' they give thanks in the Gradual (\work{De divinis officiis}~XII.23).

\subsection{The cry of those not yet gathered}
\label{sec:body-depths}

The Alleluia and the Offertory sing \latin{De profúndis clamávi ad te,
Dómine}, ``From the depths I have cried to thee, O Lord.'' St Augustine,
preaching on the psalm, recalls that Jonas cried from the whale's belly, and
says: ``we also ought to understand from what deep we cry unto the Lord. For
this mortal life is our deep.'' Those who do not cry are worse off: ``they are
very deep in the deep, who do not even cry from the deep'' (\work{Enarrationes
in Psalmos}~129.1, in the abridged English of the Nicene and Post-Nicene
Fathers). On the psalm's fourth verse, which the Mass does not sing, he says
that the ``pouring forth of innocent blood blotted out all the sins of the
guilty: so great a price paid down redeemed all captives from the hand of the
enemy who captured them'' (129.3). Since for Augustine this mortal life is
the deep, those who cry from it are those whom the angels of verse~31 have not
yet gathered.
The Offertory sings the same cry again as the gifts are brought to the altar,
and repeats its first line.

For Rupert the Offertory is the brothers' plea to Joseph: \latin{Obsecramus
ut obliviscaris sceleris fratrum}, ``I beseech thee to forget the wickedness of
thy brethren'' (Gen 50:17). The Communion is Joseph's answer, \latin{Nolite
timere. Vos cogitatis de me malum, et Deus vertit illud in bonum}, ``Fear not
\dots\ You thought evil against me: but God turned it into good'' (Gen
50:19--20).

\subsection{Hearts turned to one centre}
\label{sec:body-orations}

Jerome says that we must ``hasten to him whose lightning goes out from the
east.'' Haste of that kind starts in the will, and the Collect asks God to
``arouse \dots\ the wills of Thy faithful'' (Lasance). The Secret asks God to
``turn all our hearts unto Thee'' (Lasance). In the Latin the hearts are
\latin{ómnium nostrum}, ``of all of us,'' and they are turned \latin{ad te},
``to thee'': the prayer asks for the hearts of many to be turned to one God,
as the elect of verse~31 are brought from every quarter to one Lord.

The Postcommunion is said by those who have just received the sacraments of
the altar. Neither Jerome nor Augustine speaks of the Eucharist at this verse;
Jerome's gathering is to the Passion wherever it is read, and Augustine's
eagles are fed ``as it were'' from the body by imitating it. Dom Lucien
Fromage, continuing Dom Guéranger's \work{The Liturgical Year} in the
nineteenth century, prays on this Gospel for the elect ``whose youth, like that of the
eagle, has been so often renewed by their receiving thy precious Blood,'' and
says that nothing will stay their flight ``once the sacred prey of thy divine
Body shall be shown them,'' bringing them ``all together to the banquet of the
eternal Pasch'' (vol.~11, pp.~509--510). Fromage's eagles are renewed by
Communion, and those who say the Postcommunion have just received it. They
ask that ``whatever is evil in our hearts may be restored by its gift of
healing'' (Lasance), and they ask it still on the way to the gathering of
verse~31.

\subsection{The carcass and the exiles}
\label{sec:body-difficulty}

The difficulty is the proverb itself. Its plain picture is of birds on a dead
body, and a reader may ask whether it can bear the weight the Fathers give it.
Each Father gives his reason. Jerome's rests on the word, \latin{ptoma},
``what falls in death''; Augustine's on the speaker, who said it as one about
to die; Hilary's on the fitness of glory returning where lowliness was.
Chrysostom's rests on the context, a coming visible everywhere like lightning.
If their answers stand, the verse promises that the end will be an assembly
of the scattered to the one who died. Their difference of place remains, and
none of the four excludes the others.

A second difficulty concerns the Introit. Jeremias's promise was made to
exiles in Babylon, to be kept after seventy years (Jer 29:10), and that
remains the first sense of the letter. Jerome's anagogy takes the promise
itself further. Augustine does not comment on Jeremias: he reads the
captivity of the psalm verse as sin, and this interpretation reads the
antiphon's captivity the same way because the two captivities are sung in one
chant.

\begin{fourSenses}
\item[Literal.] Jeremias promises the exiles in Babylon that God will hear them
and bring them back from all the places where he scattered them (Jer
29:11--14). Paul thanks the Father who has delivered the Colossians from the
power of darkness into the kingdom of his Son (Col 1:12--14). Jesus tells his
disciples that after the tribulation the Son of man will come visibly, like
lightning, and that the angels will gather the elect from the four winds, as
eagles gather where the body lies (Mt 24:27--31).
\item[Allegorical.] The body is Christ crucified and risen: the Passion to
which the saints are summoned (Jerome), the place of the Passion to which the
glory returns (Hilary), the body taken into heaven (Augustine). The eagles are
the saints renewed like the eagle, with the angels and martyrs (Chrysostom).
The captivity turned back is sin forgiven (Augustine, \work{En.\ in Ps.}~84.4);
the remnant of Israel comes back to the one it rejected, as Joseph's brothers
came to him (Rupert).
\item[Moral.] To come where the body is: to gather wherever the Passion is read
(Jerome); to imitate his Passion and lowliness and so be fed from him
(Augustine); to confess the captivity one is in, ``that you may be worthy to
be set free'' (Augustine, 84.5); to cry from the deep rather than lie in it
without crying (Augustine, 129.1).
\item[Anagogical.] The last assembly of the elect from the four winds, where
the glory returns to the place of the Passion and the saints are carried over
into the honour of a spiritual substance (Hilary, 25.8, 26.1); the penitent
brought back from the confusion of this world to their own place, the Church
(Jerome, \work{In Hieremiam}~V); the banquet of the eternal Pasch
(Fromage).
\end{fourSenses}
